% /chapters/4-group-actions/02-border-tracing.tex

\section{Border tracing in oriented adjacency graphs of polygonal tilings}\label{sec:border_tracing}

The work presented in this section is a departure from the algebraic perspective of the rest of this chapter; it arose as a side project, and the full details, including proofs of all stated results, appear in~\cite{WiederholdMatosWiederhold2026}.
The connection to the present chapter is that the underlying combinatorial objects are graphs with additional structure, specifically, a cyclic ordering of the neighbours of each vertex, and the algorithms we describe amount to a systematic traversal of these structures.

\emph{Border tracing} is a classical segmentation technique in digital image processing: given a binary image, one wishes to determine a closed path of pixels that traces the boundary or contour of an object.
The standard setting is the integer lattice~$\mathbb{Z}^2$ with the well-known $4$- and $8$-adjacencies, and algorithms for this case have been described in numerous textbooks (see \cite{gonzalez2018,klette2004} for example).

A natural generalization replaces square pixels by the tiles of an arbitrary edge-to-edge polygonal tiling of $\mathbb{R}^2$, for example, a triangular or hexagonal tiling, and equips the set of tiles with an
adjacency relation.
This generalization is motivated both by theoretical considerations in digital topology and by practical applications: triangular and hexagonal pixel grids are used extensively in sensors, computer graphics, and data visualization.
In~\cite{WiederholdMatosWiederhold2026}, we develop a border tracing algorithm that works for objects in any such tiling, adapting and correcting the \emph{Border Mesh Determination} method of Voss~\cite{voss1993}.

\subsection{Oriented adjacency graphs, border meshes and Freeman codes}

Let $\mathcal{T}$ be an edge-to-edge polygonal tiling of $\mathbb{R}^2$ consisting of convex polygons, and let $A$ be an irreflexive, symmetric relation on $\mathcal{T}$ such that the graph $G = (\mathcal{T}, A)$ is connected and any two $A$-adjacent tiles intersect.
We call $A$ an \emph{adjacency} on $\mathcal{T}$; the two most natural choices are \emph{edge-adjacency} (where tiles share a side) and \emph{vertex-adjacency} (where tiles share a vertex, and possibly more).

An \emph{orientation} of $G$ assigns to each tile $t$ a cyclic ordering $z(t) = (q_1, q_2, \dots, q_{v(t)})$ of its $v(t)$ neighbours.
Starting from a directed edge $(t,s)$, one can repeatedly move to the next edge by following the neighbour of $s$ that comes immediately after $t$ in the cyclic order at $s$.
Repeating this rule produces a directed cycle or path, called a \emph{mesh}.
Because the rule is deterministic, each directed edge belongs to a unique maximal mesh.
The set of all maximal meshes, we denote by $M$.

A finite subset $\mathcal{C} \subset \mathcal{T}$ (an \emph{object}) induces a subgraph whose restricted orientation may create meshes not present in $M$.
Any such new mesh is called a \emph{border mesh} of $\mathcal{C}$.
Intuitively, border meshes trace the boundary of the object within the tiling.

When the adjacency graph is drawn in~$\mathbb{R}^2$ with each tile represented by its centroid and each edge by a straight segment, one obtains a natural \emph{counterclockwise} (respectively, \emph{clockwise}) orientation by ordering the neighbours of each tile according to the angular order of the corresponding segments.
For a connected, hole-free object under such an orientation, the border mesh found by our algorithm is called the \emph{contour}; it circumscribes the entire object in a precise sense made rigorous in \cite[Theorem~1]{WiederholdMatosWiederhold2026}.

\paragraph{Freeman code.}
When $A$ is a $k$-adjacency (every tile has exactly $k$ neighbours), the cyclic orderings can be encoded uniformly.
This is important since it turns the cyclic order into a simple integer lookup so the algorithm can move from one boundary tile to the next.
For adjacent tiles $t$ and $t'$ with $z(t) = (q_1, q_2, \dots, q_k)$ and $t' = q_i$, the \emph{Freeman code} of the ordered pair $(t,t')$ is the integer $f(t,t') = i - 1$.
Thus the cyclic order $z(t)$ is represented by the sequence $(0, 1, 2, \dots, k-1)$, and any closed path
$(t_1, t_2, \dots, t_n)$ has a \emph{Freeman chain code} $(f(t_1,t_2),\; f(t_2,t_3),\; \dots,\; f(t_{n-1},t_n),\; f(t_n,t_1))$.
This encoding is the natural extension of the classical Freeman chain code for $4$- and $8$-neighbours on the square grid, introduced in \cite{gonzalez2018}.

Figure~\ref{fig:freeman_codes} shows the Freeman codes for a triangular tiling.

\begin{figure}[ht]
    \centering
    \includegraphics[width=0.85\textwidth]{fig_freeman_codes}
    \caption{Freeman codes for the edge- and vertex-neighbours in the standard triangular tiling, providing the counterclockwise orientation.
    Reproduced from~\cite[Figure~12]{WiederholdMatosWiederhold2026}.}
    \label{fig:freeman_codes}
\end{figure}

\subsection{The algorithm BorderTracer}

The general border tracing procedure \cite[Algorithm~1]{WiederholdMatosWiederhold2026} works for any oriented adjacency graph and any object whose components each have at least two tiles.
When the adjacency is a $k$-adjacency, the algorithm simplifies to the version presented below as Algorithm~\ref{alg:border-tracing}, which uses the Freeman code to drive the neighbour search.

A key correction over the original method of Voss~\cite{voss1993} is the \textbf{End Condition Test}: a border mesh may pass through the starting tile $t_0$ multiple times before closing, and hence it is not sufficient to stop upon first return to $t_0$.
Instead, the algorithm verifies that both the first \emph{directed edge} $(t_0, t_1)$ has been repeated, that is, that $t' = t_0$, and that the subsequent tile equals $t_1$.

\begin{algorithm}[ht]
\caption{Border Tracing for a $k$-adjacency graph $G = (\mathcal{T}, A, Z, M)$}
\label{alg:border-tracing}
\DontPrintSemicolon
\SetKwInOut{Input}{Input}
\SetKwInOut{Output}{Output}
\Input{Object $\mathcal{C} \subset \mathcal{T}$;\; a finite set $\mathcal{N}$ with $\mathcal{C} \subset \mathcal{N} \subset \mathcal{T}$ surrounding the object;\; the Freeman code sequence $(0,1,\dots,k{-}1)$.}
\Output{Ordered list $\mathrm{Cont}$ of tiles forming a border mesh of~$\mathcal{C}$ in~$G$.}

\textbf{Step 1.}\;
Scan $\mathcal{N}$ to find a tile $t = t_0 \in \mathcal{C}$ having a neighbour $q \in \mathcal{N} \setminus \mathcal{C}$.\;
Set $\mathrm{Cont} \leftarrow (t_0)$.\;

\textbf{Step 2.}\;
Set $b \leftarrow f(t, q)$.\;
\Repeat{a tile $t' \in \mathcal{C}$ is found \textup{(}at most $k$ iterations\textup{)}}{
    $b \leftarrow (b+1) \bmod k$\;
    Let $t'$ be the tile with $f(t, t') = b$\;
}
\lIf{no object tile was found}{%
    \Return error: ``$\mathcal{C}$ is an isolated tile''}
Append $t'$ to $\mathrm{Cont}$, so $\mathrm{Cont} = (t_0, t_1)$.\;

\textbf{Step 3.}\;
Let $\mathrm{Cont} = (t_0, t_1, \dots, t_n)$.\;
Set $b \leftarrow f(t_n, t_{n-1})$.\;
\Repeat{a tile $t' \in \mathcal{C}$ is found}{
    $b \leftarrow (b+1) \bmod k$\;
    Let $t'$ be the tile with $f(t_n, t') = b$\;
}

\textbf{End Condition Test:}\;
\uIf{$t' \neq t_0$}{
    Append $t'$ to $\mathrm{Cont}$; go to \textbf{Step 3}.\;
}
\Else{
    Append $t'$ to $\mathrm{Cont}$.
    Run \textbf{Step 3} once more to obtain the next tile $t' \in \mathcal{C}$.\;
    \uIf{$t' = t_1$}{
    Remove the last element of $\mathrm{Cont}$.\;
    \Return $\mathrm{Cont}$.\;
    }
    \Else{
    Append $t'$ to $\mathrm{Cont}$; go to \textbf{Step 3}.\;
    }
}
\end{algorithm}

The algorithm has linear time complexity in the length of the output list $\mathrm{Cont}$: at each step, locating the previous tile in the cyclic order and scanning for the next object tile each require at most $k$ comparisons.

Figure~\ref{fig:border_mesh_example} illustrates the border meshes obtained by Algorithm~\ref{alg:border-tracing} for a connected object with a hole, using the counterclockwise orientation.

\begin{figure}[ht]
    \centering
    \includegraphics[width=0.75\textwidth]{fig_border_mesh_example}
    \caption{Border meshes of a connected object in the edge-adjacency graph of a hexagonal tiling under the counterclockwise orientation.
    Reproduced from \cite[Figure~11/13]{WiederholdMatosWiederhold2026}.}
    \label{fig:border_mesh_example}
\end{figure}

Under the counterclockwise (respectively, clockwise) orientation, and for an $A$-connected object $\mathcal{C}$ with edge-connected background, the contour produced by Algorithm~\ref{alg:border-tracing} (starting from edge\-neigh\-bours $t_0 \in \mathcal{C}$, $q \notin \mathcal{C}$) satisfies the following properties proved in~\cite{WiederholdMatosWiederhold2026}:
\begin{enumerate}
    \item \textbf{Circumscription.}
    The contour visits every tile of $\mathcal{C}$ that meets the topological frontier $\mathrm{fr}(|\mathcal{C}|)$, and the sequence of frontier segments contributed by the contour tiles traces $\mathrm{fr}(|\mathcal{C}|)$ exactly once, counterclockwise (respectively, clockwise).

    \item \textbf{Separation} (for planar edge-adjacency graphs).
    Deleting all edges between the contour and the background disconnects $\mathcal{C}$ from its complement.
    
    \item \textbf{Boundary coincidence} (under edge-adjacency).
    The set of tiles appearing in the contour equals the boundary $\mathcal{B}(\mathcal{C})$, defined as all tiles of $\mathcal{C}$ intersecting $\mathrm{fr}(|\mathcal{C}|)$.
\end{enumerate}

\paragraph{Implementation.}
A complete implementation of Algorithm~\ref{alg:border-tracing} is available in a public repository \cite{wiederhold2026border_repo}.
It supports the edge- and vertex-adjacencies on rectangular, triangular, and hexagonal tilings, and includes routines for generating the Freeman codes, detecting starting pairs~$(t_0, q)$, and extracting all border meshes via repeated application of the algorithm with appropriate marking of already-traced borders.